%

\pol{\section{Okrąg dziewięciu punktów (Feuerbacha) i prosta Eulera}}
\ita{\section{Circonferenza dei nove punti (di Feuerbach) e retta di Eulero}}

\todofoot{https://www.deltami.edu.pl/2012/09/okrag-dziewieciu-punktow-i-pewne-dwa-fakty/}
\todofoot{https://www.deltami.edu.pl/2018/04/wysokosci-czworokata/ Monge}
\todofoot{https://www.deltami.edu.pl/2020/10/jego-wysokosci-czesc-2/}
\todofoot{https://www.deltami.edu.pl/2022/05/wspolliniowosc-i-katy/ i 20-10}

\pol{Prosta Eulera to pierwsza w szkolnej geometrii trójka punktów współliniowych.}
\ita{La retta di Eulero è probabilmente il primo esempio scolastico di tre punti allineati.}
\pol{Przyszła na świat w~1765 roku, w żurnalu \emph{Novi commentarii Academiae Petropolitanae}, w artykule \emph{Solutio facilis problematum quorundam geometricorum difficillimorum}, czyli jakbyśmy powiedzieli po polsku, ,,Łatwe rozwiązanie niektórych najtrudniejszych problemów geometrycznych''.}
\ita{Vide la luce nel 1765 sulla rivista \emph{Novi commentarii Academiae Petropolitanae}, nell'articolo \emph{Solutio facilis problematum quorundam geometricorum difficillimorum}, ossia, tradotto liberamente, \emph{Facile soluzione di alcuni difficilissimi problemi geometrici}.}

\begin{proposition}[\pol{prosta Eulera}\ita{retta di Eulero}]
    \label{prosta_eulera}
    \pol{Środek okręgu opisanego na nierównobocznym trójkącie $O$, środek ciężkości $G$ oraz ortocentrum $H$ leżą na jednej prostej, zwanej prostą Eulera.}
    \ita{Il circocentro $O$, il baricentro $G$ e l'ortocentro $H$ di un triangolo scaleno giacciono su una stessa retta, detta retta di Eulero.}
    \pol{\index{prosta!Eulera}}
    \ita{\index{retta!di Eulero}}
    \pol{Mamy $|HG| = 2|OG|$.}
    \ita{Vale inoltre $|HG| = 2|OG|$.}
\end{proposition}

\pol{Piszą o niej Coxeter \cite[s. 32, 33]{coxeter_1967}, Hartshorne \cite[s. 54, 55]{hartshorne2000}, Guzicki \cite[s. 222]{guzicki_2021}, Bogdańska, Neugebauer \cite[s. 84]{neugebauer_2018}, Eves \cite[s. 109]{eves1_1972}, Audin \cite[s. 61]{audin_2003}.}
\ita{Ne parlano Coxeter \cite[p. 32, 33]{coxeter_1967}, Hartshorne \cite[p. 54, 55]{hartshorne2000}, Guzicki \cite[p. 222]{guzicki_2021}, Bogdańska e Neugebauer \cite[p. 84]{neugebauer_2018}, Eves \cite[p. 109]{eves1_1972}, Audin \cite[p. 61]{audin_2003}.}
\pol{W $\Delta_{84}^{4}$ podany będzie przepis, jak skonstruować trójkąt, w którym prosta Eulera ma zadane położenie względem podstawy.}
\ita{In $\Delta_{84}^{4}$ verrà presentato un metodo per costruire un triangolo in cui la retta di Eulero abbia una posizione assegnata rispetto alla base.}

\begin{proposition}
    \pol{Prosta Eulera jest prostopadła do jednego z boków trójkąta wtedy i tylko wtedy, gdy trójkąt jest równoramienny.}
    \ita{La retta di Eulero è perpendicolare a uno dei lati del triangolo se e solo se il triangolo è isoscele.}
\end{proposition}

\begin{proposition}
    \pol{W nierównoramiennym trójkącie, niech $d$ oznacza odległość między środkiem okręgu wpisanego i~prostą Eulera, $v$ najdłuższą środkową, $u$ najdłuższy bok, zaś $p$ połowę obwodu.}
    \ita{In un triangolo non isoscele, sia $d$ la distanza tra l'incentro e la retta di Eulero, $v$ la mediana più lunga, $u$ il lato più lungo e $p$ il semiperimetro.}
    \pol{Wtedy}
    \ita{Allora}
    \begin{equation}
        \frac ds < \frac du < \frac dv < \frac 13.
    \end{equation}
    \pol{Ograniczenia górnego nie można poprawić.}
    \ita{Il limite superiore non può essere migliorato.}
\end{proposition}
% https://en.wikipedia.org/wiki/List_of_triangle_inequalities#Euler_line

\begin{proposition}[\pol{okrąg dziewięciu punktów}\ita{circonferenza dei nove punti}]
\label{okrag_dziewieciu_punktow}%
    \pol{W każdym trójkącie środki boków, spodki wysokości oraz środki odcinków łączących ortocentrum z wierzchołkami leżą na jednym okręgu.}
    \ita{In ogni triangolo, i punti medi dei lati, i piedi delle altezze e i punti medi dei segmenti che uniscono l'ortocentro ai vertici giacciono su una stessa circonferenza.}
    \pol{Jego środek pokrywa się ze środkiem odcinka łączącego środek okręgu opisanego z ortocentrum, zaś jego promień jest dwukrotnie krótszy od promienia okręgu opisanego.}
    \ita{Il suo centro coincide con il punto medio del segmento che unisce il circocentro all'ortocentro, mentre il suo raggio è pari alla metà del raggio della circonferenza circoscritta.}
    \pol{\index{okrąg!dziewięciu punktów}}
    \ita{\index{circonferenza!dei nove punti}}
\end{proposition}

% TODO: https://en.wikipedia.org/wiki/Nine-point_circle

\pol{Coxeter \cite[s. 34, 35, 88]{coxeter_1967}, Hartshorne \cite[s. 57, 60]{hartshorne2000}, Guzicki \cite[s. 223]{guzicki_2021}, Bogdańska, Neugebauer \cite[s. 85, 86]{neugebauer_2018}, Eves \cite[s. 109]{eves1_1972}, Audin \cite[s. 62]{audin_2003}. Zetel \cite[s. 59, 85]{zetel_2020} (ze wzmianką o okręgu jedenastu punktów).}
\ita{Ne scrivono Coxeter \cite[p. 34, 35, 88]{coxeter_1967}, Hartshorne \cite[p. 57, 60]{hartshorne2000}, Guzicki \cite[p. 223]{guzicki_2021}, Bogdańska e Neugebauer \cite[p. 85, 86]{neugebauer_2018}, Eves \cite[p. 109]{eves1_1972}, Audin \cite[p. 62]{audin_2003}. Zetel \cite[p. 59, 85]{zetel_2020} la menziona assieme alla circonferenza degli undici punti.}

\pol{Temat badali Benjamin Bevan (który zasugerował środek oraz promień) i John Butterworth (który udowodnił podejrzenia Bevana) na początku XIX wieku.}
\ita{L'argomento fu studiato all'inizio del XIX secolo da Benjamin Bevan (che suggerì il centro e il raggio) e da John Butterworth (che dimostrò le congetture di Bevan).}
\index[persons]{Bevan, Benjamin}%
\index[persons]{Butterworth, John}%
\pol{To, że środki boków i spodki wysokości leżą na wspólnym okręgu, zostało zauważone w 1821 roku przez Charles Brianchona i Jean-Victora Ponceleta.}
\ita{Il fatto che i punti medi dei lati e i piedi delle altezze giacciano su una stessa circonferenza fu osservato nel 1821 da Charles Brianchon e Jean-Victor Poncelet.}
\index[persons]{Brianchon, Charles}%
\index[persons]{Poncelet, Jean-Victor}%
\pol{Tego samego odkrycia dokonał rok później Karl Feuerbach; a krótko po nim Olry Terquem zauważył, że leży na nim dziewięć, a nie tylko sześć wspomnianych punktów.}
\ita{La stessa scoperta fu fatta un anno dopo da Karl Feuerbach; poco dopo, Olry Terquem osservò che su tale circonferenza giacciono nove punti, e non soltanto sei.}
\todofoot{The nine-point circle also passes through Kimberling centers Xi for i=11 (the Feuerbach point), 113, 114, 115 (center of the Kiepert hyperbola), 116, 117, 118, 119, 120, 121, 122, 123, 124, 125 (center of the Jerabek hyperbola), 126, 127, 128, 129, 130, 131, 132, 133, 134, 135, 136, 137, 138, 139, 1312, 1313, 1560, 1566, 2039, 2040, and 2679.}
\todofoot{Karl Wilhelm Feuerbach's Eigenschaften einiger merkwiirdigen Punkte des geradlinigen Dreiecks, along with many other interesting proofs relating to the nine point circle.}
\index[persons]{Feuerbach, Karl}%
\index[persons]{Terquem, Olry}%
\pol{Terquemowi (1842) zawdzięczamy nazwę ,,okrąg dziewięciu punktów''.}
\ita{A Terquem (1842) dobbiamo il nome \emph{circonferenza dei nove punti}.}
\todofoot{The circle is officially designated the "nine point circle" (le cercle des neuf points) by Terquem, one of the editors of the Nouvelles Annales. (see Volume I page 198).}
\pol{Zetel napisze, że Euler udowodnił je w 1765 roku w Petersburgu.}
\ita{Zetel scrive che Eulero lo dimostrò già nel 1765 a San Pietroburgo.}

\pol{Feuerbach udowodnił też, że:}
\ita{Feuerbach dimostrò inoltre che:}

\begin{theorem}[\pol{Feuerbacha}\ita{di Feuerbach}]
\label{punkt_feuerbacha}%
    \pol{Okrąg dziewięciu punktów jest styczny wewnętrznie do okręgu wpisanego (w punkcie Feuerbacha) i zewnętrznie do trzech okręgów dopisanych.}
    \ita{La circonferenza dei nove punti è tangente internamente alla circonferenza inscritta (nel punto di Feuerbach) ed esternamente alle tre circonferenze exinscritte.}
    \pol{\index{twierdzenie!Feuerbacha}%
    }
    \ita{\index{teorema!di Feuerbach}%
    }
\end{theorem}

\pol{Coxeter \cite[s. 99]{coxeter_1967}, Audin \cite[s. 110]{audin_2003} podają ten fakt w formie ćwiczenia. Eves \cite[s. 133]{eves1_1972} podaje pełne uzasadnienie. Zetel \cite[s. 70]{zetel_2020} bardzo prosto dowodzi, że punkty styczności można znaleźć bez konstruowania pięciu okręgów.}
\ita{Coxeter \cite[p. 99]{coxeter_1967} e Audin \cite[p. 110]{audin_2003} presentano questo fatto come esercizio. Eves \cite[p. 133]{eves1_1972} ne fornisce una dimostrazione completa. Zetel \cite[p. 70]{zetel_2020} mostra in modo molto semplice che i punti di tangenza possono essere trovati senza costruire cinque circonferenze.}

\begin{proposition}
    \label{orthic_triangle}
    \pol{Niech $ABC$ będzie trójkątem ostrokątnym, zaś $K$, $L$ oraz $M$ spodkami jego wysokości.}
    \ita{Sia $ABC$ un triangolo acutangolo e siano $K$, $L$ e $M$ i piedi delle sue altezze.}
    \pol{Wtedy wysokości trójkąta $ABC$ są dwusiecznymi kątów trójkąta $KLM$.}
    \ita{Allora le altezze del triangolo $ABC$ sono le bisettrici degli angoli del triangolo $KLM$.}
\end{proposition}

\pol{Hartshorne \cite[s. 58]{hartshorne2000}.}
\ita{Hartshorne \cite[p. 58]{hartshorne2000}.}

\begin{theorem}[Hamiltona]
    \pol{Dany jest trójkąt $ABC$ oraz jego ortocentrum $H$.}
    \ita{Siano dati un triangolo $ABC$ e il suo ortocentro $H$.}
    \pol{Wtedy okręgi dziewięciu punktów trójkątów $ABC$ oraz $ABH$ są identyczne: mają ten sam środek i promień.}
    \ita{Allora le circonferenze dei nove punti dei triangoli $ABC$ e $ABH$ coincidono: hanno lo stesso centro e lo stesso raggio.}
\end{theorem}

\pol{O twierdzeniu Hamiltona pisze Zetel \cite[s. 65]{zetel_2020}; rzekomo pierwszy dowód poda Hamilton w 1861 roku. Bogdańska, Neugebauer \cite[s. 86]{neugebauer_2018} dodają założenie, że trójkąt nie jest prostokątny. Wynika stąd, że cztery proste Eulera czterech trójkątów przechodzą przez jeden punkt.}
\ita{Del teorema di Hamilton scrive Zetel \cite[p. 65]{zetel_2020}; a quanto pare la prima dimostrazione fu data da Hamilton nel 1861. Bogdańska e Neugebauer \cite[p. 86]{neugebauer_2018} aggiungono l'ipotesi che il triangolo non sia rettangolo. Ne segue che le quattro rette di Eulero dei quattro triangoli passano per uno stesso punto.}

\begin{corollary}
    \pol{Okręgi opisane na trójkątach $ABC$, $ABH$, $BCH$, $ACH$ mają równe promienie.}
    \ita{Le circonferenze circoscritte ai triangoli $ABC$, $ABH$, $BCH$ e $ACH$ hanno lo stesso raggio.}
\end{corollary}

%